%%%%%%%%%%%%%%%%%%%%%%%%%%%%%%%%%%%%%%%%%%%%%%%%%%%%%%%%%%%%%%%%%%%%%%%%%%%%%%%%
% Robot Learning -- course project report
%%%%%%%%%%%%%%%%%%%%%%%%%%%%%%%%%%%%%%%%%%%%%%%%%%%%%%%%%%%%%%%%%%%%%%%%%%%%%%%%

\documentclass[letterpaper, 10 pt, conference]{ieeeconf}

\IEEEoverridecommandlockouts
\overrideIEEEmargins

\usepackage{graphicx}
\usepackage{amsmath}
\usepackage{amssymb}
\usepackage{booktabs}
\usepackage{url}

\title{\LARGE \bf
From Single-Task to Multi-Task Soft Actor-Critic on Meta-World:\\
Diagnosing Why Two MT10 Tasks Never Learn
}

\author{Johannes Heidersberger and Hemjyoti Das%
\thanks{Project report for 384.195 Robot Learning, TU Wien.}%
}

\begin{document}

\maketitle
\thispagestyle{empty}
\pagestyle{empty}

%%%%%%%%%%%%%%%%%%%%%%%%%%%%%%%%%%%%%%%%%%%%%%%%%%%%%%%%%%%%%%%%%%%%%%%%%%%%%%%%
\begin{abstract}

This report documents our course project on multi-task reinforcement learning with Soft
Actor-Critic (SAC) on Meta-World, worked through in three stages: single tasks, the
three-task suite MT3, and the ten-task suite MT10. Our MT10 baseline learns eight of the ten
tasks and leaves \emph{pick-place} and \emph{peg-insert-side} at exactly zero success for as
long as we train it. Rather than tune further, we tried to establish why. Measuring the peak
per-step reward a random policy encounters gives $3.59$ on \emph{reach} but $0.03$ on
\emph{peg-insert-side}, a 120-fold spread matching the order in which the ten tasks are
learned. Since the solved tasks reach returns near $1{,}800$, the residuals belonging to the
failing tasks are negligible in the shared squared-error critic loss. Acting on this, we
normalise the critic loss by a per-task running estimate of return scale --- an adaptation
of PopArt --- and combine it with a replay curriculum weighted by measured failure rate.
Both read only statistics the run computes about itself, so no per-task constants are
involved. Mean success on MT10 rises from $0.79$ to $0.973$ at the 6M-step budget the
ablation uses, and continuing that run to 9M brings all ten tasks to at least $0.9$ with a
best evaluation of $0.993$ --- including the two that had never succeeded once under the
baseline. We report the lessons of each stage, an ablation over the two mechanisms, a
comparison against single-task SAC, and the limitations of a single-seed study.

\end{abstract}

%%%%%%%%%%%%%%%%%%%%%%%%%%%%%%%%%%%%%%%%%%%%%%%%%%%%%%%%%%%%%%%%%%%%%%%%%%%%%%%%
\section{INTRODUCTION}

Multi-task reinforcement learning asks a single policy to solve several tasks at once.
Meta-World \cite{yu2020metaworld} is the standard benchmark for this in robot manipulation:
a Sawyer arm, a shared observation and action space, and task suites of increasing size. The
assignment for this project was to train SAC \cite{haarnoja2018sac} on individual Meta-World
tasks, then on MT3, then on MT10, and to report what changes between those settings.

The multi-task agent is the natural baseline and needs no algorithmic change: all tasks run
concurrently, a task one-hot is appended to the observation, and one replay buffer is shared
--- multi-task SAC (MT-SAC). What we found is that it has a sharp, reproducible failure
mode. On MT10 eight tasks reach $1.00$ within roughly 1.5M environment steps, and
\emph{pick-place} and \emph{peg-insert-side} then sit at exactly $0.00$ indefinitely. They
do not learn slowly; they do not learn at all, and their evaluation return stays at the
reward floor from the first evaluation onwards.

The intuitive reading is a data-allocation problem --- the hard tasks are a tenth of the
experience and need a larger share --- and the common remedy is a curriculum that samples
failing tasks more often. We implemented that remedy, and it was not sufficient. Most of
this report is an account of finding out why.

What we did in this project:

\begin{itemize}
\item Trained SAC on single Meta-World tasks, MT3 and MT10, and recorded which
      configuration choices actually mattered at each stage (Sec.~\ref{sec:stages}).
\item Diagnosed the MT10 stall by measurement rather than by trial: we ruled out
      exploration collapse, replay retention and the episode-length cap with CPU probes of
      a saved policy, and identified reward scale as the binding constraint
      (Sec.~\ref{sec:mt10stall}).
\item Implemented two mechanisms in response --- a replay curriculum weighted by failure
      rate, and a per-task normalisation of the critic loss adapted from PopArt
      \cite{vanhasselt2016popart,hessel2019popart} --- and ablated them against each other
      (Sec.~\ref{sec:method}, Sec.~\ref{sec:results}).
\item Held to the constraint that no setting may name an individual task. Hand-allocating
      resources to whichever task is currently failing does work, but it amounts to tuning
      until the benchmark passes and does not carry to a larger suite.
\end{itemize}

%%%%%%%%%%%%%%%%%%%%%%%%%%%%%%%%%%%%%%%%%%%%%%%%%%%%%%%%%%%%%%%%%%%%%%%%%%%%%%%%
\section{BACKGROUND}

Everything we use is established; this section records where it comes from.

\textbf{Multi-task RL on Meta-World.} Most work on this benchmark targets representation
sharing: soft modularisation routes a shared network per task \cite{yang2020softmod}, and
CARE mixes a set of encoders using task metadata \cite{sodhani2021care}. These change what
the network computes, which is a different question from the one we ended up looking at.

\textbf{Value normalisation.} PopArt \cite{vanhasselt2016popart} normalises value targets
adaptively while preserving the network's outputs, and was applied to multi-task RL
precisely because tasks differ in reward scale \cite{hessel2019popart}. That work uses
discrete actions and separate value heads per task. Our setting has one shared
continuous-control critic whose output SAC's actor consumes directly, so we scale the loss
rather than the targets; this keeps $Q$ in real units and avoids PopArt's
output-preservation correction, at the cost of not normalising the targets themselves.

\textbf{Task curricula.} Weighting task sampling by measured failure rate is a standard
remedy for imbalanced multi-task training, typically as $p_i \propto f_i / \sum_j f_j$
applied to episode collection. We apply the weights to replay sampling instead, and add a
floor and a cap, for the reasons in Sec.~\ref{sec:curriculum}.

\textbf{Critic regularisation.} LayerNorm after each hidden critic layer is common to
several recent value-learning stabilisers \cite{nauman2024bro,lee2025simba,bhatt2024crossq}.
We adopted it from that work for the value-explosion problems described in
Sec.~\ref{sec:stages} and kept it in every run from MT3 onwards. We did not ablate it, so we
report it as an inherited design choice rather than as a finding of ours.

%%%%%%%%%%%%%%%%%%%%%%%%%%%%%%%%%%%%%%%%%%%%%%%%%%%%%%%%%%%%%%%%%%%%%%%%%%%%%%%%
\section{PRELIMINARIES AND SETUP}

We use Meta-World v3 \cite{yu2020metaworld} through Gymnasium, with Stable-Baselines3's SAC
\cite{raffin2021sb3} as the learner. For a suite of $N$ tasks, all tasks run concurrently in
a vectorised environment, an $N$-dimensional one-hot identifying the task is appended to
each observation, and one replay buffer is shared by all of them. Episodes are capped at 200
steps. Success is Meta-World's binary per-task criterion, evaluated on 30 episodes per task
every 100k environment steps.

Networks are $[512, 1024, 1024, 512]$ for both actor and critic, with LayerNorm after each
hidden critic layer. Learning rate $3\times10^{-4}$, $\tau = 10^{-3}$, $\gamma = 0.95$,
automatic entropy tuning, $\log\sigma_{\text{init}} = -3$. MT10 uses 9 environments per task
(90 total), a 3.3M-transition replay buffer, batch 384, and 60 gradient steps per collection
round, which is an update-to-data ratio of $0.67$ --- the same ratio as on MT3. All runs
were carried out on a single NVIDIA T1000 (8\,GB) with 20 CPU cores; one 6M-step MT10 run
takes approximately 32 hours.

%%%%%%%%%%%%%%%%%%%%%%%%%%%%%%%%%%%%%%%%%%%%%%%%%%%%%%%%%%%%%%%%%%%%%%%%%%%%%%%%
\section{GETTING TO MT10}
\label{sec:stages}

MT10 is the setting this report is about. The single-task and MT3 stages are included
because they are where the configuration MT10 inherits was decided, and because the MT3
result is what talked us out of per-task allocation before we ever reached MT10. Neither
stage is presented as a controlled experiment.

\subsection{Single tasks: parallelism and the discount}

Training SAC on \emph{reach}, \emph{push} and \emph{pick-place} individually gave us the
reference points in Table~\ref{tab:single} and two settings we kept for everything
afterwards.

The first is the number of parallel environments, the second the discount. Our initial
\emph{push} run used a single environment and $\gamma = 0.99$: it reached $0.20$ early, then
collapsed to $0.00$ and stayed there, reproducibly, with a diverging critic. Moving to 16
environments and $\gamma = 0.95$ took the same task to $1.00$. We changed both at once and
did not go back to separate them, so we cannot say which mattered more; what we can say is
that the combination was necessary, and we kept it for every later run.

The reasoning behind the discount is worth stating because the same quantity returns in
Sec.~\ref{sec:mt10stall}. With a per-step reward bounded by 10, the value of a solved task
is bounded by $10/(1-\gamma)$: $200$ at $0.95$ but $1{,}000$ at $0.99$. The lower ceiling
keeps the magnitudes the critic has to fit within a range where they remain comparable
across tasks.

\begin{table}[h]
\caption{Single-task SAC, 2M environment steps each, final evaluation. The last row differs
from the others in two settings at once and is reported as the failure that set the
defaults, not as an isolated effect.}
\label{tab:single}
\begin{center}
\small
\begin{tabular}{lcc}
\toprule
task & configuration & success \\
\midrule
reach            & 16 envs, $\gamma = 0.95$ & 1.00 \\
push             & 16 envs, $\gamma = 0.95$ & 0.93 \\
pick-place       & 16 envs, $\gamma = 0.95$ & 0.80 \\
push             & 1 env, $\gamma = 0.99$   & 0.00 \\
\bottomrule
\end{tabular}
\end{center}
\end{table}

\subsection{MT3: capacity, not allocation}
\label{sec:mt3}

MT3 (\emph{reach}, \emph{push}, \emph{pick-place}) is where the multi-task failure mode
first appeared, and where runs are cheap enough to iterate. Table~\ref{tab:mt3} gives three
configurations.

Our first MT3 run left \emph{pick-place} at $0.00$ after 6M steps. Hand-allocating
environments towards it ($2/4/9$ instead of $5/5/5$) fixed it, reaching $0.93$, and this is
what first suggested an allocation problem. The third row is what changed our mind:
restoring a \emph{uniform} split while scaling the learner's capacity --- more parallel
environments, a larger replay buffer, a larger batch, proportionally more gradient steps and
a LayerNorm critic --- solves all three tasks completely, and does so without allocating
anything to \emph{pick-place} in particular.

These rows are development steps rather than an ablation: each changes several settings at
once, so they cannot attribute the improvement to any single one. What they do establish is
the weaker but sufficient claim we needed --- that the hand-tuned allocation was not
necessary, because a configuration that never names a task reaches at least as far. That
was enough to justify dropping per-task allocation before moving to MT10, where a
controlled ablation is affordable because only two mechanisms are in question.

\begin{table}[h]
\caption{MT3 after 6M environment steps, final evaluation on \emph{pick-place}, the only
task of the three that is ever in doubt. These are successive development steps, not a
controlled ablation: several settings change per row. The third row additionally uses a
LayerNorm critic.}
\label{tab:mt3}
\begin{center}
\small
\begin{tabular}{lcccccc}
\toprule
configuration & envs & buffer & batch & $\gamma$ & grad. & pick-pl. \\
\midrule
initial          & 5/5/5 & 300k & 128 & 0.99 & 1  & 0.00 \\
hand allocation  & 2/4/9 & 300k & 256 & 0.95 & 10 & 0.93 \\
uniform, scaled  & 9/9/9 & 1M   & 384 & 0.95 & 18 & \textbf{1.00} \\
\bottomrule
\end{tabular}
\end{center}
\end{table}

\subsection{MT10: the stall, and what it is not}
\label{sec:mt10stall}

We carried the MT3 configuration to MT10 unchanged except for the four quantities that are
functions of task count (buffer size, batch size, gradient steps, total steps). It solved
eight of ten tasks and reproduced the stall exactly: mean success flat at $0.79$ from 1.4M
steps, with \emph{pick-place} and \emph{peg-insert-side} at $0.00$ throughout. Scaling
capacity was therefore not by itself the answer either.

Three explanations were available. We tested all three on CPU by probing the saved policy,
rather than by launching further training runs, and all three are wrong.

\emph{Exploration collapse.} If the failing tasks had lost their exploration noise they
could not discover the sparse success condition. The measured per-task action standard
deviation contradicts this as strongly as possible: \emph{peg-insert-side} is the
\emph{noisiest} task in the run at $0.877$ and \emph{pick-place} second at $0.531$, against
$0.272$ for the fully solved \emph{reach}. This also predicts that a per-task entropy
coefficient --- the standard remedy --- would push where the agent is already saturated, so
we implemented it but left it disabled.

\emph{Insufficient retention.} The shared buffer holds 3.3M transitions across ten tasks,
i.e.\ 330k per task under a uniform split, matching the retention that solved
\emph{pick-place} on MT3. Retention was already correct.

\emph{The episode cap.} Using Meta-World's scripted expert policies we measured
steps-to-success at a 500-step horizon: \emph{peg-insert-side} is solved in a median of 93
steps, and 17 of 20 attempts finish within our 200-step cap. The cap implies a ceiling near
85\% on that task, but cannot explain a floor of zero.

\subsection{What the measurement showed}

The decisive quantity turned out to be how much reward is reachable \emph{before} a task is
learned. We rolled out a random policy on each task and recorded the peak per-step reward it
encountered:

\begin{center}
\small
\begin{tabular}{lcccc}
\toprule
 & reach & push & pick-place & peg-insert-side \\
\midrule
peak reward & 3.59 & 0.54 & 0.08 & 0.03 \\
\bottomrule
\end{tabular}
\end{center}

This is a 120-fold spread, and it reproduces the order in which the ten tasks were learned.
Its consequence for a shared critic is arithmetic. Once the easy tasks are solved their
evaluation returns range from $1{,}081$ to $1{,}867$, while an unlearned
\emph{peg-insert-side} episode returns under $70$. The critic is fitted by minimising
squared error over a batch containing all tasks, so residuals enter the gradient in
proportion to the square of their magnitude: a two-order-of-magnitude gap in return scale is
a four-order-of-magnitude gap in gradient contribution. The signal belonging to
\emph{peg-insert-side} is not merely smaller than the others, it is negligible against them.
The critic reduces its loss most efficiently by ignoring the task, and the actor, which
follows $Q$, ignores it in turn.

This also explains why the curriculum was not enough. A curriculum controls \emph{how often}
a task is sampled; it cannot make what that task shows the critic register. The two are
orthogonal, and for \emph{peg-insert-side} the second is the binding constraint.

%%%%%%%%%%%%%%%%%%%%%%%%%%%%%%%%%%%%%%%%%%%%%%%%%%%%%%%%%%%%%%%%%%%%%%%%%%%%%%%%
\section{METHOD}
\label{sec:method}

\subsection{Per-task critic-loss normalisation}

We maintain, for each task $i$, an exponential moving estimate $\sigma_i$ of the standard
deviation of that task's TD targets, updated from each batch with momentum
$\alpha = 10^{-3}$:
%
$$
\sigma_i \leftarrow (1-\alpha)\,\sigma_i + \alpha \cdot \mathrm{std}\!\left(\{y_k : \tau_k = i\}\right)
\eqno{(1)}
$$
%
where $y_k$ is the target for sample $k$ and $\tau_k$ its task, recovered from the one-hot
block of the observation. The critic loss divides each sample's residual by its own task's
scale:
%
$$
\mathcal{L} = \tfrac{1}{2}\sum_{j}\; \mathbb{E}_k
\left[ \left( \frac{Q_j(s_k,a_k) - y_k}{\max(\sigma_{\tau_k},\,\sigma_{\min})} \right)^{\!2} \right]
\eqno{(2)}
$$
%
with $\sigma_{\min} = 10^{-2}$ guarding against a task whose targets are momentarily
constant. Every task now contributes comparably to the loss regardless of its reward
magnitude.

Two properties matter. We scale the \emph{loss}, not the targets, so $Q$ stays in real
units and the actor objective, the target computation and every logged value are unchanged.
And $\sigma_i$ is estimated from the run's own TD targets: it is not a constant, not a
per-task setting, and not carried over from a previous run.

\subsection{Failure-rate replay curriculum}
\label{sec:curriculum}

Let $\bar{s}_i$ be task $i$'s mean success over the last $M = 3$ evaluations. We weight
replay sampling by
%
$$
w_i \;\propto\; \mathrm{clip}\left(1 - \bar{s}_i,\; \ell,\; c\right),
\qquad \ell = 0.1,\; c = 0.6
\eqno{(3)}
$$
%
normalised over tasks and applied as a per-environment sampling probability inside the
replay buffer. Task-to-environment assignment is fixed at construction, so this costs one
weighted draw per batch and no per-transition bookkeeping.

Three choices distinguish this from the usual formulation. It acts on \textbf{replay rather
than collection}, because collection share and optimisation share are separable and only the
latter binds: under a uniform split MT10 already collects far more \emph{pick-place}
transitions than the amount that solved it standalone. The \textbf{floor} $\ell$ prevents a
solved task's share from reaching zero, which is forgetting. The \textbf{cap} $c$ prevents a
task that cannot yet be learned from absorbing the entire budget, which is what a bare
$f_i / \sum_j f_j$ rule does at ten tasks when two are stuck at zero. With eight tasks
solved and two at zero, (3) gives each failing task 115 of every 384 samples and each solved
task 19.

\subsection{No per-task constants}

No configuration file in this project names a task. The curriculum reads evaluation success,
which the run measures; the normalisation reads TD-target dispersion, which the run
computes. The only hand-chosen constants --- $\ell$, $c$, $\alpha$, $\sigma_{\min}$ and the
number of environments --- are global and shared by every task. Four further settings
(buffer size, batch size, gradient steps, total steps) follow from the number of tasks by
arithmetic. Applying the same configuration to a larger suite such as MT30 would therefore
need no further search, though we had neither the compute nor the time to verify that.

%%%%%%%%%%%%%%%%%%%%%%%%%%%%%%%%%%%%%%%%%%%%%%%%%%%%%%%%%%%%%%%%%%%%%%%%%%%%%%%%
\section{RESULTS}
\label{sec:results}

\subsection{MT10 ablation}

We ran the four combinations of the two mechanisms from scratch to 6M steps on the identical
configuration of Sec.~\ref{sec:method}; the arms differ only in which flags are set.
Fig.~\ref{fig:mean} gives mean success and Fig.~\ref{fig:hard} the two tasks the baseline
never learns; Table~\ref{tab:mt10} lists the final numbers.

\begin{figure}[t]
  \centering
  \includegraphics[width=\columnwidth]{fig/mt10_mean.pdf}
  \caption{MT10 mean success over all ten tasks, for the four combinations of replay
  curriculum and normalised critic loss. The dotted line marks the 6M-step budget at which
  the arms are compared; the combined arm alone was continued past it.}
  \label{fig:mean}
\end{figure}

\begin{figure}[t]
  \centering
  \includegraphics[width=\columnwidth]{fig/mt10_hard.pdf}
  \caption{The two tasks the baseline never learns. Both remain at exactly zero without the
  normalised critic loss.}
  \label{fig:hard}
\end{figure}

% NOTE: the curriculum-only and normalisation-only arms are still running; fill in the two
% middle rows from `python report/make_figures.py` before submitting.
\begin{table}[h]
\caption{MT10 ablation at the common 6M-step budget, final evaluation. ``Solved'' counts
tasks at or above $0.9$ out of ten. The combined arm was continued to 9M afterwards, where
it reaches $10/10$ and a best mean of $0.993$; that continuation is excluded here so the
arms remain comparable.}
\label{tab:mt10}
\begin{center}
\small
\begin{tabular}{llcccc}
\toprule
curriculum & normalisation & pick-pl. & peg-ins. & solved & mean \\
\midrule
--       & --       & 0.00 & 0.00 & 8 & 0.793 \\
\checkmark & --     & \multicolumn{4}{c}{\emph{running}} \\
--       & \checkmark & \multicolumn{4}{c}{\emph{running}} \\
\checkmark & \checkmark & \textbf{0.77} & \textbf{1.00} & 9 & \textbf{0.973} \\
\bottomrule
\end{tabular}
\end{center}
\end{table}

The baseline plateaus at $0.79$ from 1.4M steps onwards with both hard tasks at exactly
$0.00$ and their evaluation return at the reward floor.

With both mechanisms enabled, \emph{peg-insert-side} records its first non-zero evaluation
at 2.7M steps, passes $0.5$ at 3.3M and reaches $1.00$ --- all 30 evaluation episodes --- at
5.6M. \emph{Pick-place} follows the same trajectory roughly 800k steps behind. At the 6M
budget the run reaches mean success $0.973$ with nine tasks at or above $0.9$.

Because both hard tasks were still improving at 6M we continued this arm to 9M; the other
arms stop at the budget, and the tables quote the budget so the comparison stays
like-for-like. Over the last 1M steps mean success holds between $0.973$ and $0.993$,
\emph{pick-place} rises into a $0.83$--$1.00$ band and \emph{peg-insert-side} holds
$0.93$--$1.00$. Four evaluations have all ten tasks at or above $0.9$ simultaneously, the
best of them at mean $0.993$ (8.4M steps). Both tasks that the baseline left at exactly zero
are therefore solved, not merely lifted off the floor.

Two observations qualify this. First, normalisation gives no early advantage: at 2.2M the
combined arm is marginally \emph{behind} the baseline ($0.777$ against $0.793$). The
mechanism does not accelerate tasks that are already learnable; it makes tasks learnable
that otherwise never become so, and that only shows after 2.5M steps. Second, there is a
cost. \emph{Push} --- solved and stable at $1.00$ in the baseline --- becomes noisy under
the curriculum, oscillating between $0.80$ and $1.00$, because the floor $\ell = 0.1$
deliberately reduces a solved task's replay share to 19 samples per batch. The mean is far
higher regardless, but this identifies $\ell$ as the parameter to raise if a solved task
were to degrade further.

\subsection{The mechanism is visible in the statistics}

Fig.~\ref{fig:scale} plots the learned per-task value scales $\sigma_i$. They span $15$ to
$89$ at convergence, and far more early in training: \emph{peg-insert-side} begins near
$0.04$ and rises to $6.7$ as the task starts to be learned, then into the range of the
others. This is a direct measurement of the imbalance the method corrects. Without (2), a
residual on the $\sigma \approx 0.04$ task is weighted identically to one on the
$\sigma \approx 89$ task and contributes some $5\times10^{6}$ times less to the gradient.

\begin{figure}[t]
  \centering
  \includegraphics[width=\columnwidth]{fig/value_scale.pdf}
  \caption{Per-task value scale $\sigma_i$ during the combined run (log axis). The two tasks
  the baseline never learns start orders of magnitude below the rest and rise as they become
  learnable.}
  \label{fig:scale}
\end{figure}

\subsection{Multi-task against single-task}

Comparing the completed MT10 run against the single-task runs of Table~\ref{tab:single}, all
at their final evaluation: \emph{reach} $1.00$ against $1.00$, \emph{push} $1.00$ against
$0.93$, \emph{pick-place} $0.83$ against $0.80$. The multi-task agent matches or exceeds
single-task SAC on every task we trained separately, despite dividing its capacity ten ways
and using 4.5M steps per task against the single-task runs' 2M. It also solves
\emph{peg-insert-side}, which we never trained standalone. Sharing one critic across ten
tasks cost us nothing measurable here once the scale imbalance was corrected; we make no
claim of positive transfer, only that the usual multi-task penalty did not appear.

%%%%%%%%%%%%%%%%%%%%%%%%%%%%%%%%%%%%%%%%%%%%%%%%%%%%%%%%%%%%%%%%%%%%%%%%%%%%%%%%
\section{DISCUSSION AND LIMITATIONS}

\textbf{Single seed.} Every number here comes from one seed per configuration. One T1000 and
32 hours per MT10 run did not allow seed replication alongside the ablation. The central
claim is robust to this in one respect --- a task moving from exactly $0.00$ across 23
consecutive evaluations to $1.00$ on 30 episodes is not seed noise --- but the mean-success
differences between arms should be read as indicative rather than significant.

\textbf{An earlier, weaker version of the ablation.} Before running the arms from scratch we
tested the curriculum by warm-starting it from the baseline's 2.2M checkpoint. It moved
\emph{pick-place} off zero for the first time ($0.03$, $0.07$, $0.17$ over successive
evaluations) but never consolidated it, and did nothing for \emph{peg-insert-side}. That
probe is what prompted the diagnosis in Sec.~\ref{sec:mt10stall}, but it shares optimiser
state and replay contents with the baseline and covers a different step range, so we do not
report it as an arm.

\textbf{The episode cap bounds \emph{peg-insert-side}.} Our scripted-expert measurement
implies a ceiling near 85\% at 200 steps, yet the learned policy reaches $1.00$. The learned
solution is therefore faster than Meta-World's reference controller on this task, which we
did not anticipate and did not optimise for.

\textbf{Not tested.} A per-task entropy coefficient is implemented and left disabled, on the
measurement that the failing tasks were already the noisiest in the run; testing it as a
fifth arm would strengthen the comparison against the usual remedy. We also did not attempt
MT30 or MT50, so the claim that the configuration transfers by arithmetic is untested.

%%%%%%%%%%%%%%%%%%%%%%%%%%%%%%%%%%%%%%%%%%%%%%%%%%%%%%%%%%%%%%%%%%%%%%%%%%%%%%%%
\section{CONCLUSIONS}

We scaled SAC from single Meta-World tasks to MT10 and spent most of the project on the two
tasks that would not learn. The useful lesson for us was that the obvious explanation --- too
little data --- was testable and false, and that a cheap CPU measurement (the reward a random
policy can reach on each task) identified the real constraint after two full training runs
had failed to fix it. Correcting the critic loss for per-task return scale, together with a
replay curriculum, took MT10 mean success from $0.79$ to $0.973$ at a matched budget, and to
all ten tasks solved when that run was continued to 9M --- including
\emph{peg-insert-side}, which had never once succeeded beforehand.

Neither mechanism is new in itself: the normalisation is an adaptation of PopArt and the
curriculum is a standard failure-rate rule with a floor and a cap. What we would keep from
this project is the working habit --- when a run stalls, measure the quantity the algorithm
actually depends on before changing hyperparameters --- and the discipline of not letting
any setting name an individual task, which is what makes the result an argument rather than
a tuning record.

\addtolength{\textheight}{-3cm}

%%%%%%%%%%%%%%%%%%%%%%%%%%%%%%%%%%%%%%%%%%%%%%%%%%%%%%%%%%%%%%%%%%%%%%%%%%%%%%%%
\begin{thebibliography}{99}

\bibitem{yu2020metaworld} T. Yu, D. Quillen, Z. He, R. Julian, K. Hausman, C. Finn, and
S. Levine, ``Meta-World: A benchmark and evaluation for multi-task and meta reinforcement
learning,'' in \emph{Conf. on Robot Learning (CoRL)}, 2019.

\bibitem{haarnoja2018sac} T. Haarnoja, A. Zhou, P. Abbeel, and S. Levine, ``Soft
actor-critic: Off-policy maximum entropy deep reinforcement learning with a stochastic
actor,'' in \emph{Int. Conf. on Machine Learning (ICML)}, 2018.

\bibitem{vanhasselt2016popart} H. van Hasselt, A. Guez, M. Hessel, V. Mnih, and D. Silver,
``Learning values across many orders of magnitude,'' in \emph{Advances in Neural Information
Processing Systems (NeurIPS)}, 2016.

\bibitem{hessel2019popart} M. Hessel, H. Soyer, L. Espeholt, W. Czarnecki, S. Schmitt, and
H. van Hasselt, ``Multi-task deep reinforcement learning with PopArt,'' in \emph{AAAI Conf.
on Artificial Intelligence}, 2019.

\bibitem{yang2020softmod} R. Yang, H. Xu, Y. Wu, and X. Wang, ``Multi-task reinforcement
learning with soft modularization,'' in \emph{Advances in Neural Information Processing
Systems (NeurIPS)}, 2020.

\bibitem{sodhani2021care} S. Sodhani, A. Zhang, and J. Pineau, ``Multi-task reinforcement
learning with context-based representations,'' in \emph{Int. Conf. on Machine Learning
(ICML)}, 2021.

\bibitem{raffin2021sb3} A. Raffin, A. Hill, A. Gleave, A. Kanervisto, M. Ernestus, and
N. Dormann, ``Stable-Baselines3: Reliable reinforcement learning implementations,''
\emph{Journal of Machine Learning Research}, vol. 22, no. 268, pp. 1--8, 2021.

\bibitem{nauman2024bro} M. Nauman, M. Ostaszewski, K. Jankowski, P. Mi{\l}o{\'s}, and
M. Cygan, ``Bigger, regularized, optimistic: Scaling for compute and sample-efficient
continuous control,'' in \emph{Advances in Neural Information Processing Systems (NeurIPS)},
2024.

\bibitem{lee2025simba} H. Lee, D. Hwang, D. Kim, H. Kim, J. J. Tai, K. Subramanian,
P. R. Wurman, J. Choo, P. Stone, and T. Seno, ``SimBa: Simplicity bias for scaling up
parameters in deep reinforcement learning,'' in \emph{Int. Conf. on Learning Representations
(ICLR)}, 2025.

\bibitem{bhatt2024crossq} A. Bhatt, D. Palenicek, B. Belousov, M. Argus, A. Amiranashvili,
T. Brox, and J. Peters, ``CrossQ: Batch normalization in deep reinforcement learning for
greater sample efficiency and simplicity,'' in \emph{Int. Conf. on Learning Representations
(ICLR)}, 2024.

\end{thebibliography}

\end{document}
